\section{A Taxonomy of Robot-Learning Techniques}
\label{sec:taxonomy}

Figure~\ref{fig:taxonomy_main} organises the field into six branches. The dividing
question is \emph{what ships}: frozen weights (\S\ref{sec:vla}) or executable skills
(\S\ref{sec:code}). We tag mechanisms by \textbf{F}eedback, \textbf{S}earch, and
\textbf{M}emory throughout.

Table~\ref{tab:compare_main} compares all systems in the map on what each \emph{ships}, where
it is evaluated, and whether it improves from its own experience.
\input{tables/tab_compare_main}

\subsection{Survey methodology: how the corpus was built}
\label{sec:method}

Figure~\ref{fig:prisma} summarises the corpus construction in the PRISMA~2020
style~\citep{prisma2020}. The corpus was assembled in two passes over a January~2016 to
July~2026 window. First, the 77 taxonomy systems were curated by \emph{seeding and
snowballing}: for each of the six branches we started from its landmark systems and chased
citations forward and backward until new candidates stopped appearing. Second, the 225-work
landscape corpus (Table~\ref{tab:landscape}) was gathered by five \emph{structured web-search
sweeps}, one per area cluster: (i) generalist and vision-language-action policies, imitation
and diffusion policies; (ii) language-model planning, task-and-motion planning, reward and data
synthesis; (iii) skill discovery, hierarchical reinforcement learning, and world models;
(iv) manipulation, dexterity, sim-to-real, locomotion, and benchmarks; and (v) representation
learning, offline reinforcement learning, navigation, and tactile sensing. Each sweep queried
the open web with area-specific search strings of the form ``\emph{<system or family name>
robot learning arXiv}'' and ``\emph{<area> survey/benchmark/dataset}'' over arXiv, the ACM
Digital Library, IEEE Xplore, Google Scholar, and venue proceedings pages (CoRL, RSS, ICRA,
IROS, NeurIPS, ICML, ICLR, CVPR).

\emph{Inclusion} required a work to (a) concern how an embodied agent's skill or policy is
authored, learned, transferred, evaluated, or distributed, and (b) carry verifiable metadata:
exact title, first author, year, and an arXiv identifier or DOI, each confirmed against the
source. \emph{Exclusion} removed pure perception, navigation, and locomotion work with no
bearing on skill authoring (retained only in the landscape appendix where it anchors an area),
works whose metadata could not be verified, and duplicates, detected by both bibliographic key
and arXiv identifier. Every reference in this survey passed the verification step; candidates
that failed it were discarded at harvest time, and we report that their count was not logged
rather than estimate it. Included systems were then mapped onto the taxonomy of
Figure~\ref{fig:taxonomy_main} by the mechanism definitions of Table~\ref{tab:defs}.

% PRISMA-2020-style corpus-construction flow with recorded counts (see methodology subsection).
% Upright TWO-COLUMN design: two identification streams side by side, each stage's exclusions folded
% inline (red italic) rather than as separate side-boxes. Two boxes across fit the text width, so
% \resizebox scales the picture UP -> figure text renders at/above body size. Do NOT re-add separate
% side exclusion boxes (3 boxes across overflows the column and shrinks the font below body size).
\begin{figure}[t]
\centering
\resizebox{\textwidth}{!}{%
\begin{tikzpicture}[
  font=\footnotesize,
  box/.style={draw=hidden-draw, rounded corners=1pt, align=left, inner sep=5pt,
              text width=16.5em, fill=white},
  ex/.style={font=\scriptsize\itshape, text=BrickRed!85!black},
  bar/.style={draw=hidden-draw, rounded corners=1pt, inner sep=2pt, rotate=90,
              anchor=center, align=center, font=\footnotesize\bfseries, minimum width=8em},
  ar/.style={-{Stealth[length=2mm]}, darkgray, line width=0.7pt},
  node distance=4mm and 9mm,
]
% ---- left column: structured web-search stream (exclusions inline, red italic) ----
\node[box, fill=secA!30] (L1)
  {\textbf{Identification (web search).} 237 candidate records from five thematic sweeps\\
   {\scriptsize\itshape\color{BrickRed!85!black} --12 cross-sweep duplicates removed}};
\node[box, fill=secC!22, below=of L1] (L2)
  {\textbf{Screening.} 225 screened against the existing corpus (bibkey, arXiv-ID)\\
   {\scriptsize\itshape\color{BrickRed!85!black} --0 records overlapping the taxonomy corpus}};
\node[box, fill=secC!22, below=of L2] (L3)
  {\textbf{Eligibility.} 225 metadata-verified: exact title, first author, year, arXiv ID\\
   {\scriptsize\itshape\color{BrickRed!85!black} --candidates failing verification dropped at harvest (count not logged)}};
\node[box, fill=secEl, below=of L3] (L4)
  {\textbf{Landscape corpus:} 225 works (Table~\ref{tab:landscape})};
% ---- right column: seed-and-snowball curation ----
\node[box, fill=secB!25, right=9mm of L1] (R1)
  {\textbf{Identification (seeding \& snowballing).} Seed systems per branch, then citation
   chasing across the six branches};
\node[box, fill=secB!12, below=of R1] (R2)
  {\textbf{Taxonomy systems:} 77 metadata-verified};
\node[box, fill=secB!12, below=of R2] (R3)
  {\textbf{Prior surveys:} 7 (venue-targeted search)};
% ---- merged included box ----
\node[box, fill=secEl, below=7mm of L4, xshift=13em, text width=35em, align=center] (INC)
  {\textbf{Included:} 302 systems mapped (77 taxonomy $+$ 225 landscape) $+$ 7 prior surveys
   $+$ 1 industry reference $=$ \textbf{310 references}};
% ---- arrows ----
\draw[ar] (L1) -- (L2); \draw[ar] (L2) -- (L3); \draw[ar] (L3) -- (L4);
\draw[ar] (R1) -- (R2); \draw[ar] (R2) -- (R3);
\draw[ar] (L4.south) -- (L4.south |- INC.north);
\draw[ar] (R3.south) |- ([xshift=6em]INC.north) -- ([xshift=6em]INC.north |- INC.north);
% ---- stage bars ----
\node[bar, fill=secA] at ($(L1.west)+(-2em,0)$) {Identification};
\node[bar, fill=secC] at ($(L2.west)!0.5!(L3.west)+(-2em,0)$) {Screening};
\node[bar, fill=secE] at ($(L4.west)+(-2em,0)$) {Included};
\end{tikzpicture}}
\caption{\textbf{Corpus construction} in the PRISMA 2020 style~\citep{prisma2020}. The left column
is the web-search stream behind the landscape corpus (Table~\ref{tab:landscape}), with recorded
counts and inline exclusions (red) at each stage; the right column is the seed-and-snowball
curation behind the taxonomy corpus (Figure~\ref{fig:taxonomy_main}). Both streams feed the
included set. Candidates whose metadata could not be verified were discarded at harvest time; their
count was not logged, which we state rather than estimate (\S\ref{sec:method}).}
\Description{A PRISMA 2020 style flow diagram with two identification columns. The left column
(structured web search) shows 237 candidate records, minus 12 duplicates, screened to 225 with zero
taxonomy overlap, and 225 metadata-verified records forming the landscape corpus; exclusions are
noted inline, including a note that candidates failing verification were dropped with their count
not logged. The right column (seed-and-snowball curation) yields 77 taxonomy systems and 7 prior
surveys. Both columns feed an included box of 302 systems and 310 total references. Side bars mark
the Identification, Screening, and Included stages.}
\label{fig:prisma}
\end{figure}

